% Part:sets-functions-relations
% Chapter: size-of-sets
% Section: non-enumerability

\documentclass[../../../include/open-logic-section]{subfiles}

\begin{document}

\olfileid{sfr}{siz}{nen}

\olsection{\printtoken{S}{nonenumerable} ગણો}

\begin{editorial}
  આ વિભાગ \olref[enm]{sec}ની વ્યાખ્યા વાપરીને
  $\Bin^\omega$ અને $\Pow{\PosInt}$ અગણનીય છે તે સાબિત કરે છે.
  તેને \olref[enm-alt]{sec}માંના સ્વરૂપ કરતાં થોડો વધુ
  પ્રાથમિક અને થોડો વધુ વિગતવાર રહે એ રીતે રચ્યો છે.
\end{editorial}

કેટલાક ગણો, જેમ કે ધન પૂર્ણાંકોનો ગણ~$\PosInt$, અનંત છે.
અત્યાર સુધી આપણે જોયેલા અનંત ગણોના બધા દાખલા
!!{enumerable} હતા. જોકે, કેટલાક અનંત ગણોને આ ગુણધર્મ
હોતો નથી. આવા ગણોને \emph{!!{nonenumerable}} કહે છે.

સૌ પહેલાં તો !!{nonenumerable} ગણો હોય એ જ કદાચ આશ્ચર્યજનક છે.
દરેક !!{enumerable} ગણ~$A$ માટે કોઈ !!a{surjective} વિધેય
$f \colon \PosInt \to A$ હોય છે. ગણ !!{nonenumerable} હોય,
તો આવું કોઈ વિધેય હોતું નથી. એટલે કે, $\PosInt$ના અનંત ઘણા
!!{element}sનું $A$માં નિરૂપણ કરતું કોઈ વિધેય આખા~$A$ને
નિઃશેષ આવરી શકતું નથી. આમ, અનંત ઘણા ધન પૂર્ણાંકો કરતાં પણ
$A$માં જાણે ``વધુ'' !!{element}s છે.

કોઈ ગણ !!{nonenumerable} છે તે કેવી રીતે સાબિત કરવું?
આવું કોઈ વ્યાપ્ત વિધેય અસ્તિત્વમાં હોઈ શકે નહીં તે બતાવવું પડે.
સમાન રીતે કહીએ તો, $A$ના ઘટકોને એક દિશામાં અનંત સુધી જતી
યાદીમાં પરિગણિત કરી શકાતા નથી તે બતાવવું પડે. તે કરવાની ઉત્તમ
રીત એ બતાવવાની છે કે $A$ના !!{element}sની દરેક યાદીમાં ઓછામાં
ઓછો એક ઘટક રહી જ જવાનો છે; અથવા કોઈ વિધેય
$f\colon \PosInt \to A$ !!{surjective} હોઈ શકે નહીં. કૅન્ટૉરની
\emph{વિકર્ણ પદ્ધતિ}થી આ કરી શકીએ. $A$ના !!{element}sની કોઈ
યાદી, જેમ કે $x_1$, $x_2$, \dots, આપેલી હોય, તો આપણે $A$નો
એવો બીજો ઘટક બનાવીએ છીએ જે તેની રચનાને કારણે આ યાદીમાં હોઈ જ ન શકે.

આપણો પહેલો દાખલો $0$ અને $1$ના બધા અનંત, ખાલી સ્થાન વિનાના
અનુક્રમોનો ગણ~$\Bin^\omega$ છે.

\begin{thm}
\ollabel{thm:nonenum-bin-omega}
$\Bin^\omega$~એ !!{nonenumerable} છે.
\end{thm}

\begin{proof}
વિરોધાભાસ ખાતર ધારો કે $\Bin^\omega$ એ !!{enumerable} છે;
એટલે કે $\Bin^\omega$ના બધા !!{element}sની કોઈ યાદી
$s_{1}$, $s_{2}$, $s_{3}$, $s_{4}$, \dots{} છે એમ ધારો.
આ દરેક $s_i$ પોતે $0$ અને~$1$નો અનંત અનુક્રમ છે.
આ યાદીના $i$મા અનુક્રમના $j$મા ઘટકને $s_i(j)$ કહીએ.
ત્યારે $i$મો અનુક્રમ~$s_i$ આ છે:
\[
s_i(1), s_i(2), s_i(3), \dots
\]

આ યાદી અને તેમાંના દરેક અનુક્રમ~$s_i$ના ઘટકોને નીચેની
સરણિમાં ગોઠવી શકીએ:
\[
\begin{array}{c|c|c|c|c|c}
& 1 & 2 & 3 & 4 & \dots \\\hline
1 & \mathbf{s_{1}(1)} & s_{1}(2) & s_{1}(3) & s_1(4) & \dots \\\hline
2 & s_{2}(1)& \mathbf{s_{2}(2)} & s_2(3) & s_2(4) & \dots \\\hline
3 & s_{3}(1)& s_{3}(2) & \mathbf{s_3(3)} & s_3(4) & \dots \\\hline
4 & s_{4}(1)& s_{4}(2) & s_4(3) & \mathbf{s_4(4)} & \dots \\\hline
\vdots & \vdots & \vdots & \vdots & \vdots & \mathbf{\ddots}
\end{array}
\]
ડાબી બાજુનાં ચિહ્નો યાદીના અનુક્રમો $s_1$, $s_2$, \dots{}ના
ક્રમ દર્શાવે છે; ઉપરના અંકો દરેક અનુક્રમના !!{element}sને
ચિહ્નિત કરે છે. દાખલા તરીકે, $s_{1}(1)$ એ અનુક્રમ~$s_{1}$ના
પહેલા !!{element}, જે $0$ કે~$1$ હશે, તેનું નામ છે; અને એમ આગળ.

હવે આપણે $0$ અને $1$નો એવો અનંત અનુક્રમ~$\overline{s}$ બનાવીએ
જે આ યાદીમાં હોઈ જ ન શકે. $\overline{s}$ની વ્યાખ્યા યાદી
$s_1$, $s_2$, \dots પર આધાર રાખશે. $0$ અને $1$ના અનંત
અનુક્રમોની કોઈ પણ અનંત યાદીમાંથી એવો અનંત અનુક્રમ
$\overline{s}$ મળે છે જે યાદીમાં નહીં જ આવે તેની ખાતરી હોય છે.

$\overline{s}$ને વ્યાખ્યાયિત કરવા તેના બધા !!{element}s શું છે તે,
એટલે કે બધા $n \in \PosInt$ માટે $\overline{s}(n)$ શું છે તે,
નિર્દિષ્ટ કરીએ. ઉપરની સરણિના વિકર્ણને નીચેની તરફ વાંચીને
(એથી ``વિકર્ણ પદ્ધતિ'' નામ પડ્યું છે) દરેક $1$ને $0$ અને દરેક
$0$ને~$1$માં બદલીને આ કરીએ છીએ. વધુ અમૂર્ત રીતે,
વિકર્ણના $n$મા !!{element}~$s_n(n)$ની કિંમત $1$ છે કે $0$
તે પ્રમાણે $\overline{s}(n)$ને $0$ કે $1$ વ્યાખ્યાયિત કરીએ:
\[
\overline{s}(n) =
\begin{cases}
1 & \text{જો $s_{n}(n) = 0$ હોય}\\
0 & \text{જો $s_{n}(n) = 1$ હોય}.
\end{cases}
\]
કિસ્સા પાડેલી વ્યાખ્યા કરતાં સૂત્ર વધુ ગમે, તો
$\overline{s}(n) = 1 - s_n(n)$ દ્વારા પણ વ્યાખ્યાયિત કરી શકો.

સ્પષ્ટ છે કે $\overline{s}$ એ $0$ અને $1$નો અનંત અનુક્રમ છે,
કારણ કે તે આપણી સરણિના વિકર્ણ પર દેખાતા $0$ અને $1$ના
અનુક્રમને ઉલટાવીને મળતો અનુક્રમ જ છે. તેથી $\overline{s}$ એ
$\Bin^\omega$નો !!a{element} છે. પરંતુ તે યાદી
$s_1$, $s_2$, \dots{}માં હોઈ શકે નહીં. શા માટે?

તે યાદીનો પહેલો અનુક્રમ~$s_1$ હોઈ શકે નહીં, કારણ કે પહેલા
!!{element}માં તે $s_1$થી જુદો છે. $s_1(1)$ જે હોય તેનાથી
વિરુદ્ધ $\overline{s}(1)$ વ્યાખ્યાયિત કર્યું છે. તે યાદીનો
બીજો અનુક્રમ પણ હોઈ શકે નહીં, કારણ કે બીજા ઘટકમાં
$\overline{s}$ એ $s_2$થી જુદો છે: $s_2(2)$ એ $0$ હોય તો
$\overline{s}(2)$ એ $1$ છે, અને ઊલટું પણ. અને એમ આગળ.

વધુ ચોક્કસ રીતે: જો $\overline{s}$ યાદીમાં હોત, તો કોઈ $k$
માટે $\overline{s} = s_{k}$ હોત. બે અનુક્રમો ત્યારે અને માત્ર
ત્યારે જ સમાન હોય જ્યારે તેઓ દરેક સ્થાને સરખા હોય; એટલે કે કોઈ પણ
$n$ માટે $\overline{s}(n) = s_{k}(n)$. તેથી ખાસ કિસ્સા તરીકે
$n = k$ લેતાં $\overline{s}(k) = s_{k}(k)$ થવું જ પડે.
$s_k(k)$ એ $0$ કે~$1$ છે. તે $0$ હોય, તો $\overline{s}$ની
વ્યાખ્યા પ્રમાણે $\overline{s}(k)$ એ~$1$ હોવું જોઈએ. પરંતુ
$s_k(k) = 1$ હોય, તો ફરી $\overline{s}$ની વ્યાખ્યા પ્રમાણે
$\overline{s}(k) = 0$. બંને કિસ્સામાં
$\overline{s}(k) \neq s_{k}(k)$.

શરૂઆતમાં આપણે $\Bin^\omega$ના !!{element}sની યાદી
$s_1$, $s_2$, \dots{} છે એમ ધાર્યું હતું. આ યાદીમાંથી આપણે
અનુક્રમ~$\overline{s}$ બનાવ્યો અને સાબિત કર્યું કે તે યાદીમાં
હોઈ શકે નહીં. પરંતુ બધા $s_i$ એ $0$ અને $1$ના અનુક્રમો હોય,
તો તે ચોક્કસપણે $0$ અને $1$નો અનુક્રમ છે; એટલે કે
$\overline{s} \in \Bin^\omega$. ખાસ કરીને, આ બતાવે છે કે
$\Bin^\omega$ના \emph{બધા} !!{element}sની યાદી હોઈ શકે નહીં:
કોઈ પણ એવી યાદીમાંથી એવો અનુક્રમ~$\overline{s}$ બનાવી શકીએ
જે તેમાં નહીં જ આવે. તેથી $\Bin^\omega$ના બધા અનુક્રમોની યાદી
છે એવી ધારણા વિરોધાભાસ તરફ દોરી જાય છે.
\end{proof}

\begin{explain}
આ સાબિતીની રીતને ``વિકર્ણીકરણ'' કહે છે, કારણ કે તે
$\overline{s}$ની વ્યાખ્યા કરવા સરણિના વિકર્ણનો ઉપયોગ કરે છે.
વિકર્ણીકરણ માટે સરણિ હાજર જ હોય એવું જરૂરી નથી: કોઈ સરણિ કે
વાસ્તવિક વિકર્ણ ન હોય ત્યારે પણ આ જ પ્રકારનો વિચાર વાપરીને
ગણો !!{enumerable} નથી તે બતાવી શકીએ.
\end{explain}

\begin{thm}
\ollabel{thm:nonenum-pownat}
$\Pow{\PosInt}$ એ !!{enumerable} નથી.
\end{thm}

\begin{proof}
એ જ રીતે આગળ વધીએ: $\PosInt$ના ઉપગણોની દરેક યાદી માટે
$\PosInt$નો એવો ઉપગણ હોય છે જે યાદીમાં હોઈ ન શકે તે બતાવીએ.
ધારો કે નીચે $\PosInt$ના ઉપગણોની આપેલી યાદી છે:
\[
Z_{1}, Z_{2}, Z_{3}, \dots
\]
હવે ગણ~$\overline{Z}$ એવો વ્યાખ્યાયિત કરીએ કે કોઈ પણ
$n \in \PosInt$ માટે $n \in \overline{Z}$ ત્યારે અને માત્ર
ત્યારે જ જ્યારે $n \notin Z_{n}$:
\[
\overline{Z} = \Setabs{n \in \PosInt}{n \notin Z_n}
\]
$\overline{Z}$ સ્પષ્ટપણે ધન પૂર્ણાંકોનો ગણ છે, કારણ કે ધારણા
પ્રમાણે દરેક~$Z_n$ એવો છે; તેથી $\overline{Z} \in
\Pow{\PosInt}$. પરંતુ $\overline{Z}$ આ યાદીમાં હોઈ શકે નહીં.
તે બતાવવા આપણે દરેક $k \in \PosInt$ માટે $\overline{Z} \neq
Z_k$ સ્થાપિત કરીશું.

એટલે સ્વૈચ્છિક $k \in \PosInt$ લો. આપણે $\overline{Z}$ને એવો
વ્યાખ્યાયિત કર્યો છે કે કોઈ પણ $n \in \PosInt$ માટે
$n \in \overline{Z}$ ત્યારે અને માત્ર ત્યારે જ જ્યારે
$n \notin Z_n$. ખાસ કરીને, $n=k$ લેતાં
$k \in \overline{Z}$ ત્યારે અને માત્ર ત્યારે જ જ્યારે
$k \notin Z_k$. પરંતુ તેનાથી $\overline{Z} \neq Z_k$ મળે છે,
કારણ કે $k$ એકનો !!a{element} છે પણ બીજાનો નથી, અને તેથી
$\overline{Z}$ તથા $Z_k$ના !!{element}s જુદા છે. $k$ સ્વૈચ્છિક
હતો, તેથી $\overline{Z}$ યાદી $Z_1$, $Z_2$, \dots{}માં નથી.
\end{proof}

\begin{explain}
અગાઉની સાબિતીમાં વિકર્ણનો ઉલ્લેખ ન હતો, પરંતુ તેને આ રીતે
ચિત્રીને વિકર્ણ સમાયેલો હોય એમ વિચારી શકો: ગણો $Z_1$, $Z_2$,
\dots{}ને સરણિમાં લખેલા કલ્પો, જ્યાં દરેક !!{element}~$j \in Z_i$
jમા સ્તંભમાં લખાયેલો હોય. ધારો કે આ યાદીના પહેલા ચાર ગણો
$\{1,2,3,\dots\}$, $\{2, 4, 6, \dots\}$, $\{1,2,5\}$ અને
$\{3,4,5,\dots\}$ છે. ત્યારે સરણિની શરૂઆત આ રીતે થશે:
\[
\begin{array}{r@{}rrrrrrr}
  Z_1 = \{ & \mathbf{1}, & 2, & 3, & 4, & 5, & 6, & \dots\}\\
  Z_2 = \{ &  & \mathbf{2}, &  & 4, &  & 6, & \dots\}\\
  Z_3 = \{ & 1, & 2, &  &  & 5\phantom{,} &  & \}\\
  Z_4 = \{ &  &  & 3, & \mathbf{4}, & 5, & 6, & \dots\}\\
  \vdots & & & & & \ddots
\end{array}
\]
પછી વિકર્ણ પર નીચે જતાં જે અંકો દેખાય તેમને છોડી દઈને અને
$j$મી હાર તથા $j$મા સ્તંભના છેદે ખાલી જગ્યા હોય એવા~$j$નો
સમાવેશ કરીને મળતો ગણ એ $\overline{Z}$ છે. ઉપરના કિસ્સામાં આપણે
$1$ અને $2$ છોડીશું, $3$નો સમાવેશ કરીશું, $4$ છોડીશું, વગેરે.
\end{explain}

\begin{prob}
વિકર્ણ દલીલ વડે બતાવો કે $\Pow{\Nat}$ એ !!{nonenumerable} છે.
\end{prob}

\begin{prob}\label{sfr:siz:nen:prob:f-posint}
સ્પષ્ટ વિકર્ણ દલીલ વડે બતાવો કે વિધેયો
$f \colon \PosInt \to \PosInt$નો ગણ !!{nonenumerable} છે.
એટલે કે, જો $f_1$, $f_2$, \dots વિધેયોની યાદી હોય અને દરેક
$f_i\colon \PosInt \to \PosInt$ હોય, તો આ યાદીમાં ન હોય એવું
કોઈ $\overline{f}\colon \PosInt \to \PosInt$ છે તે બતાવો.
\end{prob}

\end{document}
